% ═══════════════════════════════════════════════════════════════════════════
%  SpectroClass v3.1 — Poster Científico A0
%  Congreso: Evolución estelar, exoplanetas y dinámica de sistemas estelares
%  Mendoza, 22–24 de abril de 2026
% ═══════════════════════════════════════════════════════════════════════════
\documentclass[final]{beamer}

\usepackage[size=a0,scale=1.4,orientation=portrait]{beamerposter}
\usepackage[utf8]{inputenc}
\usepackage[T1]{fontenc}
\usepackage[spanish,es-nodecimaldot]{babel}
\usepackage{lmodern}
\usepackage{amsmath,amssymb}
\usepackage{graphicx}
\usepackage{booktabs}
\usepackage{array}
\usepackage{multirow}
\usepackage{tikz}
\usepackage{tcolorbox}
\usepackage{xcolor}
\usepackage{multicol}
\usepackage{setspace}
\usepackage{microtype}
\usepackage[most]{tcolorbox}

% ── Colores institucionales ──────────────────────────────────────────────────
\definecolor{uncuyoazul}{RGB}{0,56,101}       % Azul UNCuyo
\definecolor{uncuyodorado}{RGB}{200,155,50}    % Dorado UNCuyo
\definecolor{grisclaro}{RGB}{245,245,250}
\definecolor{grismedio}{RGB}{180,180,195}
\definecolor{destacado}{RGB}{220,60,30}        % Rojo acento
\definecolor{verdeok}{RGB}{30,140,60}

% ── Tema Beamer mínimo ───────────────────────────────────────────────────────
\usetheme{default}
\usecolortheme{default}
\setbeamercolor{background canvas}{bg=white}
\setbeamercolor{block title}{fg=white,bg=uncuyoazul}
\setbeamercolor{block body}{fg=black,bg=grisclaro}
\setbeamercolor{block title alerted}{fg=white,bg=destacado}
\setbeamercolor{block body alerted}{fg=black,bg=grisclaro}
\setbeamercolor{block title example}{fg=white,bg=verdeok}
\setbeamercolor{block body example}{fg=black,bg=grisclaro}
\setbeamerfont{block title}{size=\Large,series=\bfseries}
\setbeamerfont{block body}{size=\normalsize}

% ── Caja de resultado destacado ──────────────────────────────────────────────
\tcbset{
  resultado/.style={
    enhanced,colback=uncuyodorado!15,colframe=uncuyodorado,
    fonttitle=\bfseries\large, title=#1, arc=4pt,
    boxrule=1.5pt, left=6pt, right=6pt, top=4pt, bottom=4pt
  },
  metodo/.style={
    enhanced, colback=uncuyoazul!8, colframe=uncuyoazul!50,
    fonttitle=\bfseries\normalsize, title=#1, arc=3pt,
    boxrule=1pt, left=5pt, right=5pt, top=3pt, bottom=3pt
  },
  stat/.style={
    enhanced,colback=white,colframe=uncuyoazul,
    arc=6pt, boxrule=2pt, halign=center,
    fontupper=\bfseries\huge\color{uncuyoazul},
    fontlower=\small\color{black},
    left=4pt, right=4pt, top=6pt, bottom=4pt,
    attach boxed title to top center={yshift=-3mm},
    boxed title style={colback=uncuyoazul,colframe=uncuyoazul,arc=3pt}
  }
}

% ── Sin navegación ───────────────────────────────────────────────────────────
\setbeamertemplate{navigation symbols}{}
\setbeamertemplate{headline}{}
\setbeamertemplate{footline}{}

% ─────────────────────────────────────────────────────────────────────────────
\begin{document}
\begin{frame}[t]{}

%% ════════════════ CABECERA ══════════════════════════════════════════════════
\begin{tikzpicture}[remember picture, overlay]
  %% Fondo cabecera
  \fill[uncuyoazul] (0,0) rectangle (\paperwidth, 14cm);
  %% Franja dorada inferior
  \fill[uncuyodorado] (0,0) rectangle (\paperwidth, 1.2cm);
\end{tikzpicture}

\vspace{-0.5cm}
\begin{columns}[T]
  \begin{column}{0.82\paperwidth}
    \color{white}
    {\fontsize{60}{68}\bfseries\selectfont SpectroClass v3.1\par}
    \vspace{0.4cm}
    {\fontsize{36}{42}\selectfont
     Sistema Automatizado de Clasificación Espectral Estelar\\
     con Arquitectura Multi-Método\par}
    \vspace{0.6cm}
    {\fontsize{26}{30}\selectfont
     \textbf{Roberto Butrón} \hspace{1em}
     Facultad de Ciencias Exactas y Naturales, UNCuyo · Mendoza, Argentina\par}
    \vspace{0.3cm}
    {\fontsize{22}{26}\selectfont\itshape
     Congreso: Evolución estelar, exoplanetas y dinámica de sistemas estelares
     · 22--24 de abril de 2026\par}
  \end{column}
  \begin{column}{0.15\paperwidth}
    %% Espacio para logo institucional (descomentar si se dispone de archivo)
    % \includegraphics[width=0.12\paperwidth]{logo_uncuyo.png}
    \vspace{3cm}
    \begin{center}
      \tikz\draw[white,line width=3pt,rounded corners=8pt]
        (0,0) rectangle (7cm,7cm)
        node[midway,white,font=\footnotesize\bfseries,align=center]
        {Logo\\UNCuyo};
    \end{center}
  \end{column}
\end{columns}

\vspace{0.8cm}

%% ════════════════ CUERPO — 3 COLUMNAS ═══════════════════════════════════════
\begin{columns}[t,totalwidth=\paperwidth]

%%────────────────────────────────────────────────────────────────────────────
%% COLUMNA 1
%%────────────────────────────────────────────────────────────────────────────
\begin{column}{0.32\paperwidth}

  %% ─── Motivación ────────────────────────────────────────────────────────
  \begin{block}{1 · Motivación}
    \begin{itemize}\setlength\itemsep{6pt}
      \item Los \textit{surveys} modernos (SDSS, LAMOST, Gaia) generan
            \textbf{millones} de espectros que no pueden clasificarse manualmente.
      \item La clasificación MK tradicional es subjetiva y no trazable.
      \item Los sistemas existentes priorizan velocidad o exactitud,
            pero raramente ambas \emph{con} justificación física.
      \item No existe un sistema de código abierto con documentación en
            \textbf{castellano} que integre métodos físicos y aprendizaje
            automático.
    \end{itemize}
  \end{block}

  \vspace{0.5cm}

  %% ─── Objetivo ──────────────────────────────────────────────────────────
  \begin{alertblock}{2 · Objetivo}
    Desarrollar un sistema que clasifique espectros estelares en la secuencia
    OBAFGKM + clase de luminosidad I--V, combinando métodos físicos y de ML
    con \textbf{trazabilidad diagnóstica completa} y distribución como
    aplicación web de código abierto.
  \end{alertblock}

  \vspace{0.5cm}

  %% ─── Pipeline ──────────────────────────────────────────────────────────
  \begin{block}{3 · Pipeline de Pre-procesamiento}

    \begin{tcolorbox}[metodo={Paso 1 — Rechazo de rayos cósmicos (MAD)}]
      \textbf{Sigma-clipping iterativo} usando la Desviación Absoluta de la
      Mediana ($\hat\sigma = 1.4826\,\mathrm{MAD}$) en lugar de $\sigma$
      clásico. Más robusto ante distribuciones asimétricas y picos CCD.
    \end{tcolorbox}

    \vspace{0.35cm}

    \begin{tcolorbox}[metodo={Paso 2 — Normalización al continuo (2 pasos)}]
      \textbf{Primer pase:} percentil P90 en ventanas de 50 px + spline
      cúbico.\\[4pt]
      \textbf{Segundo pase:} rechazo relativo al continuo \emph{local}
      (no global) → evita distorsión en espectros con fuerte gradiente de
      color (tipos O--B y M).
    \end{tcolorbox}

    \vspace{0.35cm}

    \begin{tcolorbox}[metodo={Paso 3 — Medición de EW (85 líneas)}]
      Integración trapezoidal $\mathrm{EW} = \int(1-F_\lambda/F_c)\,d\lambda$.\\[4pt]
      Retorna tupla $(EW,\,\mathrm{depth},\,\mathrm{quality\_flag},\,
      \lambda_\mathrm{med})$.\\[4pt]
      Ventanas adaptativas: $w \approx 4\times\mathrm{FWHM}$ para líneas
      metálicas.
    \end{tcolorbox}

    \vspace{0.35cm}

    \begin{tcolorbox}[metodo={Paso 4 — 21 Razones Diagnóstico}]
      Cubren la secuencia completa: ionización O--B
      (He\,\textsc{i}/He\,\textsc{ii}, Si\,\textsc{iii}/Si\,\textsc{ii},
      N\,\textsc{v}/N\,\textsc{iii}), temperatura F--K (Cr\,\textsc{i}/
      Fe\,\textsc{i}), luminosidad G--M (Sr\,\textsc{ii}/Fe\,\textsc{i},
      TiO/CaH).
    \end{tcolorbox}

  \end{block}

  \vspace{0.5cm}

  %% ─── Datos de entrada ─────────────────────────────────────────────────
  \begin{block}{4 · Datos de Entrada}
    \centering
    \begin{tabular}{lll}
      \toprule
      \textbf{Formato} & \textbf{Resolución} & \textbf{Rango} \\
      \midrule
      TXT (2 col.) & Cualquiera $\geq 2000$ & 3900--6800~Å \\
      FITS (WCS) & $R \sim 42\,000$ & Flexible \\
      \bottomrule
    \end{tabular}\\[0.4cm]
    Catálogo de validación: \textbf{ELODIE} (Haute-Provence)\\
    856 espectros, $R\sim42\,000$, S/N $>100$
  \end{block}

\end{column}

%%────────────────────────────────────────────────────────────────────────────
%% COLUMNA 2
%%────────────────────────────────────────────────────────────────────────────
\begin{column}{0.34\paperwidth}

  %% ─── Arquitectura multi-método ─────────────────────────────────────────
  \begin{block}{5 · Arquitectura Multi-Método}

  \begin{center}
  \begin{tikzpicture}[
    node distance=0.7cm,
    box/.style={draw=uncuyoazul,fill=uncuyoazul!10,rounded corners=4pt,
                minimum width=6.8cm,minimum height=1.1cm,
                font=\bfseries\normalsize,align=center},
    vote/.style={draw=uncuyodorado!80,fill=uncuyodorado!20,rounded corners=4pt,
                minimum width=9cm,minimum height=1.3cm,
                font=\bfseries\large,align=center},
    arr/.style={->,>=stealth,line width=1.5pt,color=uncuyoazul!70}
  ]
    \node[vote] (data) {Espectro de entrada (.txt / .fits)};

    \node[box,below=0.6cm of data,xshift=-5.2cm] (norm)
      {Normalización\\continuo (MAD)};
    \node[box,below=0.6cm of data,xshift=5.2cm] (ew)
      {Medición EW\\85 líneas};

    \node[box,below=1.5cm of data,xshift=-7.5cm] (fis)
      {\textbf{Físico}\\28 nodos MK};
    \node[box,below=1.5cm of data,xshift=-2.5cm] (rf)
      {\textbf{RandomForest}\\24 features};
    \node[box,below=1.5cm of data,xshift=2.5cm] (tmpl)
      {\textbf{Template}\\$\chi^2$ reducido};
    \node[box,below=1.5cm of data,xshift=7.5cm] (nn)
      {\textbf{k-NN / CNN-1D}\\espectro remuestreado};

    \node[vote,below=3.2cm of data] (voto)
      {Votación Ponderada ($w_\mathrm{DT}=0.4,\;w_\mathrm{neural}=0.4,\ldots$)};

    \node[vote,below=0.7cm of voto,fill=uncuyoazul!25,draw=uncuyoazul] (out)
      {Tipo MK + Clase Luminosidad I--V + Confianza \% + Alternativas};

    \draw[arr] (data) -- (norm);
    \draw[arr] (data) -- (ew);
    \draw[arr] (norm) -- (fis);
    \draw[arr] (norm) -- (rf);
    \draw[arr] (ew) -- (tmpl);
    \draw[arr] (ew) -- (nn);
    \draw[arr] (fis) -- (voto);
    \draw[arr] (rf) -- (voto);
    \draw[arr] (tmpl) -- (voto);
    \draw[arr] (nn) -- (voto);
    \draw[arr] (voto) -- (out);
  \end{tikzpicture}
  \end{center}

  \end{block}

  \vspace{0.4cm}

  %% ─── Árbol físico 28 nodos ─────────────────────────────────────────────
  \begin{block}{6 · Árbol Físico de 28 Nodos (Gray \& Corbally)}
    \begin{columns}[t]
      \begin{column}{0.48\linewidth}
        \textbf{10 nodos nuevos en v3.1:}
        \begin{itemize}\setlength\itemsep{4pt}
          \item O3--O4: N\,\textsc{v} + O\,\textsc{v}~5114
          \item O tardíos: N\,\textsc{iii}~4634--4641
          \item B con O\,\textsc{ii}
          \item F5--F9: Ca\,\textsc{i}/H$\delta$ + Sr\,\textsc{ii}~4077
          \item F/G con banda CH\,G
        \end{itemize}
      \end{column}
      \begin{column}{0.48\linewidth}
        \begin{itemize}\setlength\itemsep{4pt}
          \item G: Fe\,\textsc{i}~4144/H$\delta$ + Y\,\textsc{ii}~4376
          \item G: tripleta Mg~$b$
          \item K: Cr\,\textsc{i} 4254/4275/4290~Å (termómetro)
          \item K: Na\,\textsc{i}~D
          \item M tardías: VO + CaH
        \end{itemize}
      \end{column}
    \end{columns}
    \vspace{0.3cm}
    Cada decisión cita la línea, el umbral de EW y el ratio que la motivó
    → \textbf{trazabilidad diagnóstica completa}.
  \end{block}

  \vspace{0.4cm}

  %% ─── Resultados clave ──────────────────────────────────────────────────
  \begin{block}{7 · Resultados Principales}

    \vspace{0.3cm}
    \begin{columns}[c]
      \begin{column}{0.32\linewidth}
        \begin{tcolorbox}[stat,title={}]
          86.5\%\tcblower Exactitud global\\(votación ponderada)
        \end{tcolorbox}
      \end{column}
      \begin{column}{0.32\linewidth}
        \begin{tcolorbox}[stat,title={}]
          +2.3 pp\tcblower Mejora sobre mejor\\método individual
        \end{tcolorbox}
      \end{column}
      \begin{column}{0.32\linewidth}
        \begin{tcolorbox}[stat,title={}]
          856\tcblower Espectros ELODIE\\validados
        \end{tcolorbox}
      \end{column}
    \end{columns}

    \vspace{0.5cm}

    \begin{center}
    \renewcommand{\arraystretch}{1.3}
    \begin{tabular}{lcc}
      \toprule
      \textbf{Método} & \textbf{Exactitud} & \textbf{CV 5-fold} \\
      \midrule
      RandomForest    & 84.2\,\%  & $82.5\pm3.1$\,\% \\
      k-NN ($k=5$)    & 81.3\,\%  & $79.8\pm2.8$\,\% \\
      MKCLASS (ref.)  & $\sim$85\,\% & — \\
      StarNet (ref.)  & 89\,\%    & — \\
      \midrule
      \textbf{SpectroClass v3.1} & \textbf{86.5\,\%} & — \\
      \bottomrule
    \end{tabular}
    \end{center}
  \end{block}

\end{column}

%%────────────────────────────────────────────────────────────────────────────
%% COLUMNA 3
%%────────────────────────────────────────────────────────────────────────────
\begin{column}{0.32\paperwidth}

  %% ─── Matriz de confusión ───────────────────────────────────────────────
  \begin{block}{8 · Matriz de Confusión (N\,=\,171)}
    \centering
    \renewcommand{\arraystretch}{1.35}
    \begin{tabular}{c|ccccccc}
      \toprule
      \small\diagbox{Real}{Pred.} &
      \textbf{O} & \textbf{B} & \textbf{A} &
      \textbf{F} & \textbf{G} & \textbf{K} & \textbf{M} \\
      \midrule
      O & \cellcolor{verdeok!30}\textbf{2} & 1 & 0 & 0 & 0 & 0 & 0 \\
      B & 0 & \cellcolor{verdeok!30}\textbf{9} & 3 & 1 & 0 & 0 & 0 \\
      A & 0 & 2 & \cellcolor{verdeok!30}\textbf{24} & 2 & 0 & 0 & 0 \\
      F & 0 & 0 & 2 & \cellcolor{verdeok!30}\textbf{63} & 0 & 5 & 0 \\
      K & 0 & 0 & 0 & 5 & 0 & \cellcolor{verdeok!30}\textbf{46} & 0 \\
      M & 0 & 0 & 0 & 0 & 0 & 1 & \cellcolor{verdeok!30}\textbf{5} \\
      \bottomrule
    \end{tabular}\\[0.4cm]
    \small
    Confusión B$\leftrightarrow$A: gradiente continuo de temperatura.\\
    Confusión F$\leftrightarrow$K: solapamiento de líneas metálicas.
  \end{block}

  \vspace{0.5cm}

  %% ─── Líneas diagnóstico ────────────────────────────────────────────────
  \begin{block}{9 · 85 Líneas Diagnóstico}
    \centering
    \renewcommand{\arraystretch}{1.2}
    \begin{tabular}{lrl}
      \toprule
      \textbf{Especie} & \textbf{N} & \textbf{Uso} \\
      \midrule
      Balmer H\,\textsc{i} & 5 & Temperatura A--F \\
      He\,\textsc{i} + He\,\textsc{ii} & 10 & Tipos O--B \\
      N, C, O ionizados & 13 & Subtipos O, B \\
      Silicio (Si) & 8 & Subtipos B \\
      Calcio (Ca) & 6 & F--G--K \\
      Hierro (Fe) & 13 & G--K--M \\
      Bandas TiO & 6 & Tipo M \\
      VO, CaH & 4 & M tardías \\
      Otros metales & 20 & F--K, luminosidad \\
      \midrule
      \textbf{Total} & \textbf{85} & \\
      \bottomrule
    \end{tabular}
  \end{block}

  \vspace{0.5cm}

  %% ─── Clase de Luminosidad ──────────────────────────────────────────────
  \begin{block}{10 · Clase de Luminosidad I--V}
    \begin{itemize}\setlength\itemsep{5pt}
      \item \textbf{Ia/Ib} Supergigantes:
            Sr\,\textsc{ii}/Fe\,\textsc{i} $\gg 1$,
            líneas de Balmer estrechas.
      \item \textbf{II--III} Gigantes:
            Ba\,\textsc{ii}~4554 / Fe\,\textsc{i} elevado.
      \item \textbf{IV} Subgigantes:
            Ca\,\textsc{i}/Fe\,\textsc{i} intermedio.
      \item \textbf{V} Enanas:
            H$\beta$ ancho (ensanchamiento Stark),
            TiO/CaH $\approx 1$ en M enanas.
    \end{itemize}
    \vspace{0.3cm}
    Salida: \texttt{G2V}, \texttt{K3\,III}, \texttt{B5\,Ia}, etc.
  \end{block}

  \vspace{0.5cm}

  %% ─── Interfaz web ──────────────────────────────────────────────────────
  \begin{block}{11 · Interfaz Web Flask}
    \begin{itemize}\setlength\itemsep{5pt}
      \item Subida de espectros TXT / FITS.
      \item Visualización interactiva con líneas diagnóstico marcadas.
      \item Árbol de clasificación visual navegable.
      \item Procesamiento por lotes (CSV).
      \item Entrenamiento y métricas de modelos neuronales.
      \item Exportación CSV / TXT / PDF.
      \item Documentación integrada en \textbf{castellano}.
    \end{itemize}
  \end{block}

  \vspace{0.5cm}

  %% ─── Conclusiones ──────────────────────────────────────────────────────
  \begin{exampleblock}{12 · Conclusiones}
    \begin{enumerate}\setlength\itemsep{5pt}
      \item SpectroClass v3.1 alcanza \textbf{86.5\,\%} de exactitud sobre
            ELODIE, superando en $+2.3$~pp al mejor método individual.
      \item La normalización MAD de 2 pasos reduce errores sistemáticos en
            espectros de alta resolución y tipos O--B tempranos.
      \item El árbol de 28 nodos proporciona clasificación trazable y
            verificable para cada espectro.
      \item La votación ponderada combina robustez física con potencia de ML.
      \item Sistema distribuido como aplicación web de código abierto con
            documentación completa en castellano.
    \end{enumerate}
  \end{exampleblock}

  \vspace{0.5cm}

  %% ─── Referencias ───────────────────────────────────────────────────────
  \begin{block}{Referencias}
    \small
    \begin{itemize}\setlength\itemsep{2pt}
      \item Gray \& Corbally (2009). \textit{Stellar Spectral Classification}.
            Princeton Univ.\ Press.
      \item Gray \& Corbally (2014). \textit{AJ}, 147, 80.
      \item Luo et al.\ (2015). \textit{RAA}, 15, 1095 (LAMOST DR1).
      \item Fabbro et al.\ (2018). \textit{MNRAS}, 475, 2978 (StarNet).
      \item Moultaka et al.\ (2004). \textit{PASP}, 116, 693 (ELODIE).
    \end{itemize}
  \end{block}

\end{column}
\end{columns}

%% ════════════════ PIE DE PÁGINA ═════════════════════════════════════════════
\vspace{0.4cm}
\begin{tikzpicture}[remember picture,overlay]
  \fill[uncuyoazul] (0,0) rectangle (\paperwidth,2.4cm);
  \fill[uncuyodorado] (0,2.4cm) rectangle (\paperwidth,2.8cm);
\end{tikzpicture}

\begin{center}
  \color{white}
  {\large\bfseries Roberto Butrón \enspace·\enspace
  Facultad de Ciencias Exactas y Naturales, UNCuyo \enspace·\enspace
  Mendoza, Argentina \enspace·\enspace Abril 2026}\\[0.2cm]
  {\normalsize Código abierto \enspace·\enspace
  Documentación en castellano \enspace·\enspace
  Catálogo de validación: ELODIE (Haute-Provence)}
\end{center}

\end{frame}
\end{document}
